% META: {"id": "TNSI-BDD-EVAL-B-corrige", "chapitre": "TNSI-BASES-DE-DONNEES", "type_objet": "correction", "evaluation_ref": "TNSI-BDD-EVAL-B", "capacites": ["T-BDD-02", "T-BDD-03D", "T-BDD-03E", "T-BDD-03G", "T-BDD-03H"], "mode_creation": "ex_nihilo", "sources_inspiration": [], "status": "approved"}
% BEGIN-VERIFY
% import sqlite3
% conn = sqlite3.connect(":memory:")
% conn.execute("CREATE TABLE Employe(id INTEGER PRIMARY KEY, nom TEXT, id_service INTEGER)")
% conn.execute("CREATE TABLE Service(id INTEGER PRIMARY KEY, nom TEXT)")
% conn.executemany("INSERT INTO Service VALUES (?,?)", [(1, "Comptabilite"), (2, "Informatique")])
% conn.executemany("INSERT INTO Employe VALUES (?,?,?)", [(1, "Ben Ali", 2), (2, "Chraibi", 1), (3, "Haddad", 2)])
% conn.commit()
% conn.execute("UPDATE Employe SET id_service = 1 WHERE nom = 'Ben Ali'")
% conn.execute("DELETE FROM Employe WHERE id_service = 1")
% conn.commit()
% cur = conn.execute("SELECT nom FROM Employe")
% assert cur.fetchall() == [("Haddad",)]
% END-VERIFY
\begin{corrige}{TNSI-BDD-EVAL-B-EX1}
\textbf{Q1.} Persistance, concurrence, efficacité, sécurisation.

\textbf{Q2.} Un fichier CSV n'offre aucune gestion de la \textbf{concurrence} : si deux guichets modifient le fichier en même temps, l'une des deux réservations peut être écrasée ou perdue. Un SGBD sérialise les accès concurrents pour garantir qu'aucune réservation n'est perdue.
\end{corrige}

\begin{corrige}{TNSI-BDD-EVAL-B-EX2}
\textbf{Q1.}
\begin{sql}
SELECT Employe.nom, Service.nom
FROM Employe
JOIN Service ON Employe.id_service = Service.id
ORDER BY Employe.nom ASC;
\end{sql}

\textbf{Q2.}
\begin{sql}
UPDATE Employe SET id_service = 1 WHERE nom = 'Ben Ali';
\end{sql}

\textbf{Q3.}
\begin{sql}
DELETE FROM Employe WHERE id_service = 1;
\end{sql}
Après la question précédente, \texttt{'Ben Ali'} et \texttt{'Chraibi'} sont tous deux dans le service $1$ : il ne reste que \texttt{'Haddad'}. Sans clause \texttt{WHERE}, \texttt{DELETE FROM Employe;} supprimerait \textbf{tous} les employés, quel que soit leur service.
\end{corrige}
